\paragraph{アーキテクチャと設計の関係}
アーキテクチャと設計の境界は、しばしば曖昧である。アーキテクチャはソフトウェア設計から発展してきた分野であり、設計の一部として扱われることもある。しかし、実務上は区別して考えると理解しやすい。

\par\smallskip\noindent\refstepcounter{table}\label{tab:ch02-05}\textbf{表 \thetable: 第02章 ソフトウェアアーキテクチャ 表05}\par\smallskip
\begin{longtable}{P{0.22\linewidth}P{0.38\linewidth}P{0.32\linewidth}}
\toprule
観点 & アーキテクチャ & 設計 \\
\midrule
扱う範囲 & システム全体、環境、主要構成、品質特性、進化原則である。 & 構成要素の詳細、データ構造、アルゴリズム、実装方法である。 \\
要求との関係 & 要求を受けるだけでなく、利害関係者との交渉を通じて要求を形作ることがある。 & 既に定義された要求とアーキテクチャ制約を満たす方法を具体化する。 \\
変更影響 & 後から変更すると広範囲に影響しやすい。 & 比較的局所的に変更できる場合が多い。 \\
主な成果物 & アーキテクチャ記述、ビュー、意思決定記録、原則である。 & 詳細設計、インタフェース仕様、クラス設計、アルゴリズム設計である。 \\
\bottomrule
\end{longtable}

\begin{statementbox}{見分け方}
「この判断を後から変えると、多くの構成要素、品質特性、組織分担、運用方法に影響するか」と問う。影響が広ければ、アーキテクチャ判断である可能性が高い。
\end{statementbox}
